\section{Conclusion}\label{sec:conclusion110}
Neural Bellman Operators concern continuation learning and feasible decisions under the original recursive economic objective. This revision connects their construction to action-aligned tail accuracy, original-model Bellman structure, completed-incumbent improvement and paid stopping. The remaining-gain rule has an explicit economic meaning: omitted computation can sacrifice at most a verified amount of original-policy improvement. Its recursive extension requires a comparator-future residual and risk-distorted propagation, not an unqualified average of physical state frequencies.

The original investment comparison executes that trade-off with common information and strong conventional controls. It records substantial savings in support and certification work, strictly improved fitted policies relative to completed plain incumbents, and the small but sometimes positive cost of stopping earlier than the full safeguard. Exact ReLU--hinge policy equivalence and the competitiveness of quadratic, spline and semi-analytic methods remain visible. No internal fitting metric or favorable timing subset is substituted for original economic loss.

The research program continues to target neural-specific representation and learning advantages under equal information and resources, prospective validation of tail-aware acquisition, genuine longer-horizon and multiple-control continuous-risk comparisons, and independently calibrated economic adoption. The recursive-utility BSDE bridge makes one rigorous connection to NDU/FBSDE research without transferring results across objectives or information structures. Current successful computations and complete proofs do not close those broader claims, nor do the preserved negative results replace their pursuit. The original NBO identity and the economic and numerical-methods ambition are retained.
